Panel B reveals a clear hierarchy of vulnerability across eight sociodemographic and behavioral subgroups. Groups that are traditionally marginalized or biologically high risk face the greatest structural barriers to timely vaccination. Specifically, respondents with neutral institutional trust (65.88\% attack rate), primary education (65.75\%), and age 65+ (65.38\%) show infection risks substantially above the population average. 

The elevated $\kappa$ among adults aged 65+ (0.327) and those with medical conditions (0.288) suggests a ``double burden'' of risk: biological vulnerability is compounded by behavioral resistance to administrative delay, creating a feedback loop that degrades system resilience. These findings indicate that a ``one-size-fits-all'' vaccination strategy will systematically fail the most vulnerable segments of society. To enhance resilience, policy interventions must move beyond aggregate supply metrics and focus on reducing the specific behavioral ``time-tax'' faced by high-$\kappa$ subgroups, particularly through targeted trust-building and streamlined delivery for the elderly and those with chronic health needs.
